\section{Introduction}
\label{sec:introduction}


Neural scaling laws show that performance improves predictably with model size, dataset scale, and compute, pushing stochastic optimization toward increasingly long training runs under tight efficiency constraints~\cite{hestness2017deep, kaplan2020scaling, henighan2020scaling, hoffmann2022chinchilla}. In this regime, optimization stability becomes a practical requirement, since unstable training can slow convergence and produce critical loss spikes~\cite{bai2026spikes}. Stability is usually maintained through heuristics such as learning-rate decay, gradient clipping, and, more recently, tuning of Adam's second-moment hyperparameter~\cite{loshchilov2016sgdr,pascanu2013difficulty,bai2026spikes}. These heuristics can suppress spikes~\cite{zhang2022opt}, but oscillation can also coexist with decreasing training loss over longer intervals~\cite{cohen2021edge,damian2023self}, so a smoother loss curve need not mean faster loss reduction. The central question is therefore: which changes should a learning intervention target to improve training?

The standard explanation is geometric. Training becomes unstable when updates are too large relative to sharp directions of the loss landscape, so that they overshoot rather than converge~\cite{polyak1964methods, bertsekas1997nonlinear}. The break-even framework adds gradient noise, arguing that sharpness and gradient variance jointly shape the trajectory, especially early in training~\cite{jastrzebski2020break}. Edge-of-Stability studies show that deep networks often keep making progress near or beyond the fixed-quadratic sharpness threshold~\cite{cohen2021edge, arora2022understanding, damian2023self}, and Adam's preconditioner can produce loss spikes even on a quadratic objective~\cite{kingma2014adam,bai2026spikes}. These accounts describe whether a step overshoots the loss landscape, a landscape that is complex, problem-specific, and tied to architecture~\cite{ly2025multifractal,yang2021taxonomizing,ma2022multiscale,li2018visualizing}. They do not say how much of the motion later updates undo, nor where in the network the reversal occurs; Central Flows explicitly leaves open how loss dynamics relate to the roles of individual layers~\cite{cohen2025central}.

We take a complementary view based on how updates interact. Consecutive updates either reinforce each other, adding to net displacement, or oppose each other, undoing part of the preceding change. We call persistent reinforcement a constructive regime and persistent opposition a destructive regime. Both can be measured from consecutive gradients already computed during training, without Hessian evaluations, and the same measurement applies to the hidden representations of each layer. This lets us locate where reversal occurs in the network and test whether limiting representation displacement, rather than the size of the whole update, improves training.

In summary, we make four main contributions:
\begin{itemize}
\item \textbf{Constructive and destructive regimes.} Starting from the SGD update, we express net parameter displacement through gradient magnitudes and alignment, and define the regimes by the sign of consecutive alignment \(\rho_t\). Across MLP, ResNet, and Transformer, destructive regimes show back-and-forth motion while training loss continues to decrease; in the highest-learning-rate GD trajectories, net displacement retains only about \(2.6\)--\(5.4\%\) of the distance traveled. In continuations from trained states, smaller batches attenuate consecutive alignment toward zero, weakening the oscillation.

\item \textbf{Curvature and loss-increase diagnostics from training gradients.} Consecutive full gradients and the applied update recover directional curvature averaged over the completed step, exactly for smooth, fixed objectives and independently of the optimizer, without additional backward passes or Hessian--vector products. For full-batch GD on a quadratic objective, a completed step increases the loss precisely when \(\rho_t<-\|g_{t-1}\|/\|g_t\|\), that is, when reversal is combined with sufficient gradient growth. Under Adam, this condition predicts next-step loss increases more accurately than rules based on the previous loss rise, reversal alone, or norm growth alone.

\item \textbf{Functional tracing.} We propose a tracing algorithm that follows representation changes on fixed inputs from the outputs toward earlier layers, measuring reinforcement or reversal, magnitude, accumulated displacement, and shared timing. The resulting oscillations are approximately low-dimensional, and net output displacement retains only about \(1.7\)--\(3.5\%\) of the distance traveled. Their location differs across architectures: reversal is homogeneous across the MLP, whereas earlier ResNet and Transformer representations can remain constructive while their outputs reverse.

\item \textbf{Representation displacement as an intervention target.} We introduce a distributed displacement cap that scales down Adam's proposed update, unit by unit, according to the direct response change it would produce. In matched continuations from identical training states, we compare it with unchanged Adam, learning-rate reduction, second-moment changes, and gradient clipping. Across MLP, ResNet, and Transformer, every tested cap level reduces mean training loss and accumulated upward loss spikes relative to unchanged Adam, although functional oscillation persists.

\end{itemize}

The remainder of the paper is organized as follows. Section~\ref{sec:background} positions these contributions in related work. Sections~\ref{sec:displacement}--\ref{sec:gd-experiments} develop and test the GD/SGD framework; Sections~\ref{sec:adam}--\ref{sec:vc-interventions} connect Adam diagnostics to functional measurements and interventions. Section~\ref{sec:conclusion} discusses implications, limitations, and future work.
